\documentclass[11pt,a4paper]{beamer}
\usetheme{Madrid}
\usepackage[utf8]{inputenc}
\usepackage{amsmath}
\usepackage{amsfonts}
\usepackage{amssymb}
\usepackage{graphicx}
\usepackage{tikz}
\usepackage{booktabs}
\usepackage{array}
\usepackage{colortbl}
\usepackage{multicol}

% ============================================================
% PROFESSIONAL FONT SETTINGS
% ============================================================
\usefonttheme{professionalfonts}

\setbeamerfont{normal text}{size=\small}
\setbeamerfont{itemize/enumerate body}{size=\small}
\setbeamerfont{itemize/enumerate subbody}{size=\footnotesize}
\setbeamerfont{block body}{size=\small}
\setbeamerfont{block title}{size=\small}
\setbeamerfont{frametitle}{size=\large}
\setbeamerfont{title}{size=\large}
\setbeamerfont{subtitle}{size=\normalsize}
\setbeamerfont{author}{size=\footnotesize}
\setbeamerfont{date}{size=\footnotesize}

\newcommand{\redflag}[1]{\textcolor{red}{\textbf{[!] #1}}}
\newcommand{\greennote}[1]{\textcolor{green!50!black}{\textbf{[CORRECT] #1}}}
\newcommand{\examfocus}[1]{\textcolor{orange!80!black}{\textbf{[FOCUS] #1}}}
\newcommand{\trap}[1]{\textcolor{purple}{\textbf{[TRAP] #1}}}
\newcommand{\correct}[1]{\textcolor{green!50!black}{\textbf{[right] #1}}}
\newcommand{\scenario}[1]{\textcolor{blue!70!black}{\textbf{[SCENARIO] #1}}}
\newcommand{\keyterm}[1]{\textcolor{red!80!black}{\textbf{[KEY] #1}}}

\title[CAMS Exam Preparation]{CAMS Exam Preparation}
\subtitle{Bitcoin and Cryptocurrency Questions}
\author{All 4 Domains}
\date{}

\setbeamercolor{title}{fg=white,bg=blue!70!black}
\setbeamercolor{frametitle}{fg=white,bg=blue!70!black}
\setbeamertemplate{navigation symbols}{}

\begin{document}

% TITLE SLIDE
\begin{frame}
\titlepage
\centering
\footnotesize 40 Practice Questions | Money Laundering | VASPs | Blockchain | Regulations
\end{frame}

% DOMAIN A - QUESTION 1
\begin{frame}{Domain A: Money Laundering Methods}
\scenario{A criminal earns \$2 million from darknet drug sales in Bitcoin. To hide the origin, they split the BTC into 200 small transactions, send them through 3 mixing services, then consolidate into new wallets before transferring to an unregulated VASP for cash conversion.}
\vspace{0.2cm}
\textbf{Which money laundering techniques are being used?}
\vspace{0.2cm}
\begin{itemize}
    \item A) Structuring and trade-based laundering
    \item B) Microstructuring and mixing/tumbling
    \item C) Smurfing and cash smuggling
    \item D) Shell company formation only
\end{itemize}
\vspace{0.2cm}
\trap{This is just regular structuring.}
\correct{Answer: B - Microstructuring (small digital transactions) AND mixing/tumbling.}
\redflag{Microstructuring is specifically used with digital asset laundering. Mixers make transaction tracing extremely difficult.}
\end{frame}

% DOMAIN A - QUESTION 2
\begin{frame}{Domain A: Privacy Coins}
\scenario{A customer receives monthly payments in Monero (XMR), a privacy coin. The bank cannot trace the origin of these funds. The customer refuses to provide source of funds documentation. Declared annual income is \$60,000 but Monero deposits total \$20,000 per month.}
\vspace{0.2cm}
\textbf{How should the bank classify the money laundering risk?}
\vspace{0.2cm}
\begin{itemize}
    \item A) Lower risk than Bitcoin because Monero has stronger security
    \item B) Same risk as Bitcoin
    \item C) Heightened risk due to anonymity features
    \item D) No risk assessment needed
\end{itemize}
\vspace{0.2cm}
\trap{All crypto has same risk level.}
\correct{Answer: C - Privacy coins with non-public blockchains represent HIGHER risk. Income inconsistency requires EDD.}
\end{frame}

% DOMAIN A - QUESTION 3
\begin{frame}{Domain A: Stablecoin Risks}
\scenario{A customer purchases \$500,000 in USDT from an unregulated VASP, transfers to multiple wallets, then converts to cash with minimal KYC.}
\vspace{0.2cm}
\textbf{What is the primary money laundering vulnerability of fiat-collateralized stablecoins?}
\vspace{0.2cm}
\begin{itemize}
    \item A) They cannot be traced on blockchain
    \item B) Ease of conversion back to fiat combined with regulatory gaps
    \item C) They are always fully backed by audited reserves
    \item D) They cannot be used for cross-border transfers
\end{itemize}
\vspace{0.2cm}
\trap{Stablecoins are fully backed and low risk.}
\correct{Answer: B - Ease of conversion back into traditional fiat and regulatory gaps make stablecoins attractive for money laundering.}
\end{frame}

% DOMAIN A - QUESTION 4
\begin{frame}{Domain A: NFT Wash Trading}
\scenario{A criminal creates 20 NFT artworks, lists them for 100 ETH each, then uses 10 shell wallets to purchase the NFTs from themselves using illicit Bitcoin. Total value laundered: \$3 million.}
\vspace{0.2cm}
\textbf{What money laundering technique does this represent?}
\vspace{0.2cm}
\begin{itemize}
    \item A) Trade-based money laundering
    \item B) NFT wash trading / self-laundering
    \item C) Traditional structuring
    \item D) Cash smuggling
\end{itemize}
\vspace{0.2cm}
\trap{This is legitimate art trading.}
\correct{Answer: B - NFT wash trading. Red flags include overpricing, same entity buying and selling, and anonymous platforms.}
\end{frame}

% DOMAIN A - QUESTION 5
\begin{frame}{Domain A: Crypto Mixers}
\scenario{A VASP receives multiple deposits from wallets that have passed through ChipMixer, Blender.io, and Wasabi Wallet's CoinJoin. The customer cannot explain why mixing services were used.}
\vspace{0.2cm}
\textbf{What is the VASP's obligation?}
\vspace{0.2cm}
\begin{itemize}
    \item A) No action; mixing is legal
    \item B) File SAR and conduct EDD
    \item C) Only report if mixer is sanctioned
    \item D) Close account immediately
\end{itemize}
\vspace{0.2cm}
\trap{Only sanctioned mixers are problematic.}
\correct{Answer: B - When mixing services are used, fund tracing becomes extremely difficult. Any funds from known mixers require EDD and potential SAR.}
\end{frame}

% DOMAIN A - QUESTION 6
\begin{frame}{Domain A: CBDC AML Advantages}
\scenario{A central bank considers launching a CBDC to replace 40 percent of cash in circulation. Current tax evasion through undeclared cash is \$15 billion annually.}
\vspace{0.2cm}
\textbf{What AML advantage does a CBDC have over physical cash?}
\vspace{0.2cm}
\begin{itemize}
    \item A) CBDC cannot be used for any illegal activity
    \item B) Real-time monitoring and greater transaction transparency
    \item C) CBDC has no AML advantages
    \item D) CBDC eliminates all financial crime
\end{itemize}
\vspace{0.2cm}
\trap{Cash and CBDC are the same.}
\correct{Answer: B - Unlike cash, CBDCs can be monitored in real time, providing authorities with greater transparency to combat money laundering, tax evasion, and terrorist financing.}
\end{frame}

% DOMAIN A - QUESTION 7
\begin{frame}{Domain A: Terrorist Financing via Crypto}
\scenario{A charity front organization collects Bitcoin and Monero donations from sympathizers in 20 countries, moves funds through mixers over 12 months, converts to USDT through unregulated VASPs, and sends to wallets in a conflict zone. Total: \$5 million. The charity reports only \$200,000 in traceable donations.}
\vspace{0.2cm}
\textbf{What indicates terrorist financing rather than traditional money laundering?}
\vspace{0.2cm}
\begin{itemize}
    \item A) Use of cryptocurrency only
    \item B) High transaction volumes
    \item C) Legitimate-looking charities, small transactions, mix of funds, and informal systems
    \item D) Large cash transactions only
\end{itemize}
\vspace{0.2cm}
\trap{Only illegal funds are used for terrorist financing.}
\correct{Answer: C - Terrorist financing differs: source may be legitimate, funds move through charities, transaction sizes are often small to avoid detection.}
\end{frame}

% DOMAIN A - QUESTION 8
\begin{frame}{Domain A: DeFi Vulnerabilities}
\scenario{A DeFi protocol allows lending, borrowing, and swapping across 12 blockchains with no KYC. No identity verification required. \$80 million in ransomware proceeds has been laundered through the protocol.}
\vspace{0.2cm}
\textbf{What are the primary AML vulnerabilities?}
\vspace{0.2cm}
\begin{itemize}
    \item A) No vulnerabilities; DeFi is compliant by design
    \item B) Anonymity, cross-chain bridging, no KYC
    \item C) Only the lending feature
    \item D) Only the swapping feature
\end{itemize}
\vspace{0.2cm}
\trap{DeFi is decentralized; no one is responsible.}
\correct{Answer: B - DeFi vulnerabilities include anonymity, cross-chain transfers, and no KYC. Persons with control or sufficient influence over DeFi protocols may be considered VASPs.}
\end{frame}

% DOMAIN A - QUESTION 9
\begin{frame}{Domain A: Crypto Mining AML}
\scenario{A commercial mining operation with 10,000 ASIC miners mines 30 BTC monthly. It sells BTC to three VASPs and through P2P exchanges with no KYC. The operation claims no AML obligations because mining is not a financial service.}
\vspace{0.2cm}
\textbf{Does the mining operation have AML obligations?}
\vspace{0.2cm}
\begin{itemize}
    \item A) No; mining is purely computational
    \item B) Yes; when converting BTC to fiat, miner may need MSB/VASP registration
    \item C) Only if mining illegally
    \item D) Only for transaction fees, not block rewards
\end{itemize}
\vspace{0.2cm}
\trap{Mining is industrial, not financial.}
\correct{Answer: B - When selling BTC for fiat, the miner is engaging in money transmission. Many jurisdictions require MSB/VASP registration for commercial mining operations.}
\end{frame}

% DOMAIN A - QUESTION 10
\begin{frame}{Domain A: Algorithmic Stablecoin Exploitation}
\scenario{A criminal organization purchases \$5 million of an algorithmic stablecoin at \$1.00 using illicit funds. They trigger a collapse to \$0.10, then claim a \$4.5 million capital loss on taxes while earning \$6 million in short profits. Net: \$5 million illicit in, \$6.5 million clean out.}
\vspace{0.2cm}
\textbf{What technique does this represent?}
\vspace{0.2cm}
\begin{itemize}
    \item A) Trade-based money laundering
    \item B) Loss harvesting disguised as trading losses, plus market manipulation
    \item C) Traditional structuring
    \item D) Cash smuggling
\end{itemize}
\vspace{0.2cm}
\trap{This is just a bad investment.}
\correct{Answer: B - Algorithmic stablecoins' volatility can be exploited for loss harvesting and tax fraud. Illicit funds become legitimate trading losses for tax purposes.}
\end{frame}

% DOMAIN B - QUESTION 11
\begin{frame}{Domain B: VASP Definition}
\scenario{A company offers crypto-to-crypto exchange (BTC/ETH, no fiat), non-custodial wallets, and trading pairs. No fiat currency on-ramp. Users identified only by email. Annual volume: \$1 billion.}
\vspace{0.2cm}
\textbf{Under FATF standards, is this company a VASP?}
\vspace{0.2cm}
\begin{itemize}
    \item A) No; only companies handling fiat are VASPs
    \item B) Yes; crypto-to-crypto exchanges are VASPs regardless of fiat
    \item C) Only if holding customer funds
    \item D) Only if volume exceeds \$2 billion
\end{itemize}
\vspace{0.2cm}
\trap{No fiat means no regulation.}
\correct{Answer: B - VASPs include exchanges between virtual assets. Lack of customer identification is a red flag.}
\end{frame}

% DOMAIN B - QUESTION 12
\begin{frame}{Domain B: FATF Travel Rule for Crypto}
\scenario{A VASP receives a \$30,000 USDC transfer from another VASP. The incoming message contains only a username CryptoUser123 and a wallet address. No legal name, no address.}
\vspace{0.2cm}
\textbf{Under the FATF Travel Rule, is this compliant?}
\vspace{0.2cm}
\begin{itemize}
    \item A) Yes; transaction hash is sufficient
    \item B) No; missing required originator information
    \item C) Travel Rule does not apply to crypto
    \item D) Only need originator name; address optional
\end{itemize}
\vspace{0.2cm}
\trap{Travel Rule only applies to banks.}
\correct{Answer: B - Travel Rule applies to VASPs same as banks. Required: originator name, address, account number. Username is insufficient.}
\end{frame}

% DOMAIN B - QUESTION 13
\begin{frame}{Domain B: DeFi VASP Status}
\scenario{A DeFi protocol has governance tokens: Founder A (40 percent), Founder B (35 percent), Founder C (15 percent), Public (10 percent). Multi-sig requires 3 of 5 signatures, including both founders. Protocol collects \$10 million annual fees. No KYC.}
\vspace{0.2cm}
\textbf{Who is likely considered a VASP under FATF guidance?}
\vspace{0.2cm}
\begin{itemize}
    \item A) No one; DeFi is decentralized
    \item B) Only the public token holders
    \item C) Founders A and B (control via tokens and multi-sig)
    \item D) Only users of the protocol
\end{itemize}
\vspace{0.2cm}
\trap{DeFi equals no responsible party.}
\correct{Answer: C - FATF guidance: persons with control or sufficient influence are VASPs. Indicators: governance token ownership, multi-sig control, fee collection.}
\end{frame}

% DOMAIN B - QUESTION 14
\begin{frame}{Domain B: Blockchain Analytics}
\scenario{A ransomware victim paid 150 BTC. The investigator traces through 200 addresses and finds one address linked to a regulated VASP through KYC data.}
\vspace{0.2cm}
\textbf{How can the investigator identify the suspect?}
\vspace{0.2cm}
\begin{itemize}
    \item A) Tracing beyond 200 addresses is impossible
    \item B) On-chain analytics plus off-chain data from VASP KYC
    \item C) Wait for suspect to confess
    \item D) Only on-chain analytics work
\end{itemize}
\vspace{0.2cm}
\trap{Bitcoin is anonymous; cannot trace.}
\correct{Answer: B - On-chain analytics trace transactions. Off-chain analytics integrate VASP KYC data, social media, and other sources to link addresses to real-world identities.}
\end{frame}

% DOMAIN B - QUESTION 15
\begin{frame}{Domain B: Crypto ATM AML}
\scenario{A crypto ATM operator has 1,000 machines. Daily volume: \$5 million. No KYC for transactions under \$1,000. No MSB registration. No SAR filings. 45 percent of transactions link to known scam wallets.}
\vspace{0.2cm}
\textbf{Does the operator have AML obligations?}
\vspace{0.2cm}
\begin{itemize}
    \item A) No; ATMs are automated machines
    \item B) Yes; operator is a VASP/MSB with full AML obligations
    \item C) Only for transactions over \$10,000
    \item D) Only if operating an exchange
\end{itemize}
\vspace{0.2cm}
\trap{Crypto ATMs are unregulated.}
\correct{Answer: B - Crypto ATMs are VASPs/MSBs. Obligations: registration, KYC/CDD, transaction monitoring, SAR filing.}
\end{frame}

% DOMAIN B - QUESTION 16
\begin{frame}{Domain B: Lightning Network Challenges}
\scenario{A criminal opens a Lightning channel with \$1 million, then routes 100,000 small payments through 2,000 nodes over 6 months. Only channel opening and closing appear on-chain.}
\vspace{0.2cm}
\textbf{What AML challenge does the Lightning Network present?}
\vspace{0.2cm}
\begin{itemize}
    \item A) No challenges; Lightning is transparent
    \item B) Off-chain payments leave no public record; cannot be traced
    \item C) Only large payments are hidden
    \item D) All payments appear on-chain within 24 hours
\end{itemize}
\vspace{0.2cm}
\trap{Lightning is just faster Bitcoin.}
\correct{Answer: B - Millions of off-chain Lightning payments are not recorded on the public blockchain. Only channel opening and closing are visible to investigators.}
\end{frame}

% DOMAIN B - QUESTION 17
\begin{frame}{Domain B: Non-Custodial Wallet Risk}
\scenario{A bank customer uses MetaMask wallet. Receives crypto from 200 addresses, sends to 500 addresses monthly. Transfers fiat from VASP to bank account. Cannot explain source of crypto. Refuses to provide transaction history.}
\vspace{0.2cm}
\textbf{What is the bank's obligation?}
\vspace{0.2cm}
\begin{itemize}
    \item A) No obligation; wallet is outside bank's scope
    \item B) EDD on source of funds; file SAR if customer refuses
    \item C) Ignore; crypto is unregulated
    \item D) Close account immediately
\end{itemize}
\vspace{0.2cm}
\trap{Bank only responsible for fiat.}
\correct{Answer: B - Bank has fiat relationship. Customer's crypto activity is relevant to source of funds. Refusal to provide information is a red flag requiring SAR.}
\end{frame}

% DOMAIN B - QUESTION 18
\begin{frame}{Domain B: Cross-Chain Bridges}
\scenario{A criminal moves funds: BTC to renBTC on Ethereum to USDC on Solana to BNB on BSC to AVAX on Avalanche to XMR. Total bridges: 7. Time: 4 hours. Tracing cost: \$500,000.}
\vspace{0.2cm}
\textbf{What vulnerability do cross-chain bridges present?}
\vspace{0.2cm}
\begin{itemize}
    \item A) Bridges are fully transparent
    \item B) Bridges break transaction trails across blockchains
    \item C) Only stablecoin bridges are problematic
    \item D) All bridges require KYC
\end{itemize}
\vspace{0.2cm}
\trap{Bridges are traceable like normal transactions.}
\correct{Answer: B - Each bridge crossing creates a break in the forensic trail. Multiple bridges across multiple chains exponentially increases tracing difficulty.}
\end{frame}

% DOMAIN B - QUESTION 19
\begin{frame}{Domain B: Staking and Yield Farming ML}
\scenario{A criminal deposits 2,000 ETH (\$4 million) from phishing scams into a DeFi staking pool with no KYC. Earns 6 percent APY for 12 months. Withdraws original plus rewards. Claims funds are legitimate DeFi investment returns.}
\vspace{0.2cm}
\textbf{How does staking facilitate money laundering?}
\vspace{0.2cm}
\begin{itemize}
    \item A) Staking is fully transparent
    \item B) Creates legitimate-looking returns, provides time-based layering
    \item C) Staking requires KYC
    \item D) No laundering possible; funds locked
\end{itemize}
\vspace{0.2cm}
\trap{Staking is legitimate investment.}
\correct{Answer: B - Staking provides legitimate-looking investment returns as explanation, creates time-based layering, and generates plausible documentation to support false source of funds claims.}
\end{frame}

% DOMAIN B - QUESTION 20
\begin{frame}{Domain B: VASP Licensing Jurisdiction}
\scenario{A VASP incorporated in Country A (no crypto laws), servers in Country B, management in Countries C, D, and E, customers in 60 countries. Claims registered nowhere. Annual volume: \$10 billion.}
\vspace{0.2cm}
\textbf{Which jurisdictions have regulatory responsibility?}
\vspace{0.2cm}
\begin{itemize}
    \item A) No jurisdiction
    \item B) Country A (incorporation) and countries with meaningful presence
    \item C) Only Country B
    \item D) Only customers' home countries
\end{itemize}
\vspace{0.2cm}
\trap{No physical presence equals no regulation.}
\correct{Answer: B - FATF expects licensing where incorporated and where management located. VASPs with no clear jurisdiction are high risk.}
\end{frame}

% DOMAIN C - QUESTION 21
\begin{frame}{Domain C: MiCA Regulation}
\scenario{A crypto exchange wants to offer services in all 27 EU member states. Currently registered only in Estonia.}
\vspace{0.2cm}
\textbf{Under MiCA, what is required?}
\vspace{0.2cm}
\begin{itemize}
    \item A) Register separately in each of 27 countries
    \item B) Single MiCA license with passporting rights to all EU states
    \item C) No license required
    \item D) Only register with ESMA directly
\end{itemize}
\vspace{0.2cm}
\trap{Each country has separate crypto rules.}
\correct{Answer: B - MiCA provides single EU-wide licensing. Authorized in one member state equals passporting rights to all EU member states.}
\end{frame}

% DOMAIN C - QUESTION 22
\begin{frame}{Domain C: MiCA Stablecoin Rules}
\scenario{A company issues a stablecoin (ART) pegged to fiat and commodities. Market cap: €8 billion. Issuer in non-EU country. No EU license. Wants to serve EU customers.}
\vspace{0.2cm}
\textbf{What are the requirements under MiCA?}
\vspace{0.2cm}
\begin{itemize}
    \item A) No requirements; stablecoins exempt
    \item B) Must obtain EU license, maintain liquid reserves, comply with ART authorization
    \item C) Only need registration
    \item D) Stablecoins over €1 billion prohibited
\end{itemize}
\vspace{0.2cm}
\trap{Stablecoins are unregulated.}
\correct{Answer: B - ART issuers must be authorized in an EU member state, maintain sufficient liquid reserves (1:1 backing), and have qualified directors with good repute.}
\end{frame}

% DOMAIN C - QUESTION 23
\begin{frame}{Domain C: EU 5AMLD Prepaid Card Threshold}
\scenario{A financial institution offers prepaid cards: €200 (cash, no ID), €400 (email only), €900 (basic ID), €250 (crypto purchase). All exceed €150.}
\vspace{0.2cm}
\textbf{Under 5AMLD, which cards are non-compliant?}
\vspace{0.2cm}
\begin{itemize}
    \item A) All are compliant
    \item B) All cards exceeding €150 require KYC
    \item C) Only €900 card
    \item D) Only cards purchased with crypto
\end{itemize}
\vspace{0.2cm}
\trap{4AMLD threshold was €250, so €200 is fine.}
\correct{Answer: B - 5AMLD lowered threshold from €250 to €150. Any card exceeding €150 requires customer identification.}
\end{frame}

% DOMAIN C - QUESTION 24
\begin{frame}{Domain C: US AML Act 2020}
\scenario{A US crypto exchange with \$20 billion annual volume has no MSB registration, no AML program, no SAR filings. Claims not a financial institution.}
\vspace{0.2cm}
\textbf{Is the exchange subject to BSA requirements?}
\vspace{0.2cm}
\begin{itemize}
    \item A) No; crypto exchanges exempt
    \item B) Yes; exchanges are MSBs with full BSA requirements
    \item C) Only for fiat transactions
    \item D) Only for derivatives trading
\end{itemize}
\vspace{0.2cm}
\trap{Crypto is unregulated.}
\correct{Answer: B - AML Act 2020 includes cryptocurrency exchanges as MSBs with same licensing and reporting requirements as traditional financial institutions.}
\end{frame}

% DOMAIN C - QUESTION 25
\begin{frame}{Domain C: OFAC Sanctions and Mixers}
\scenario{A US VASP identifies transactions involving Tornado Cash (sanctioned mixer), North Korean hacking group wallet (sanctioned), and Wasabi Wallet (not sanctioned).}
\vspace{0.2cm}
\textbf{Which require blocking under OFAC?}
\vspace{0.2cm}
\begin{itemize}
    \item A) None
    \item B) Tornado Cash and North Korean group only
    \item C) All three
    \item D) Only Tornado Cash
\end{itemize}
\vspace{0.2cm}
\trap{Only sanctioned entities, not mixers.}
\correct{Answer: B - OFAC applies to designated persons and entities. Tornado Cash sanctioned; North Korean group sanctioned; Wasabi Wallet not specifically sanctioned (requires EDD).}
\end{frame}

% DOMAIN C - QUESTION 26
\begin{frame}{Domain C: FATF Recommendation 15}
\scenario{A bank develops crypto custody service. Launch planned for Month 7. AML risk assessment scheduled for Month 8. CCO informed at Month 6.}
\vspace{0.2cm}
\textbf{Is this compliant with FATF Recommendation 15?}
\vspace{0.2cm}
\begin{itemize}
    \item A) Yes; assessment after launch is acceptable
    \item B) No; assessment must be completed BEFORE launch
    \item C) Only needed if volume exceeds \$1 million
    \item D) CCO notification is sufficient
\end{itemize}
\vspace{0.2cm}
\trap{Assessment can be done anytime.}
\correct{Answer: B - FATF Recommendation 15 requires identification and assessment of ML/TF risks BEFORE launching new products or technologies.}
\end{frame}

% DOMAIN C - QUESTION 27
\begin{frame}{Domain C: EU AMLA Supervisory Powers}
\scenario{AMLA identifies a cross-border crypto exchange with €100 billion volume, customers in 20 EU states, 4 compliance failures, conflicting national supervision.}
\vspace{0.2cm}
\textbf{What powers does AMLA have?}
\vspace{0.2cm}
\begin{itemize}
    \item A) No powers; AMLA only coordinates
    \item B) Direct supervision, on-site inspections, corrective measures, penalties
    \item C) Only can issue guidance
    \item D) Only can supervise banks, not crypto
\end{itemize}
\vspace{0.2cm}
\trap{AMLA only coordinates.}
\correct{Answer: B - AMLA has direct supervision authority over high-risk cross-border entities including VASPs. Powers include on-site inspections, corrective measures, and penalties.}
\end{frame}

% DOMAIN C - QUESTION 28
\begin{frame}{Domain C: FATF Grey List Impact}
\scenario{A jurisdiction on FATF Grey List has no VASP licensing, no Travel Rule, no crypto SAR filing. 100 plus VASPs operate from there. Crypto volume: \$200 billion annually.}
\vspace{0.2cm}
\textbf{What actions should financial institutions take?}
\vspace{0.2cm}
\begin{itemize}
    \item A) No action; Grey List is advisory
    \item B) Apply enhanced due diligence to Grey List jurisdictions
    \item C) Terminate all relationships immediately
    \item D) Only apply EDD if on Black List
\end{itemize}
\vspace{0.2cm}
\trap{Grey List equals no action needed.}
\correct{Answer: B - FATF calls for enhanced due diligence for Grey List jurisdictions. Countermeasures may be required for persistent deficiencies.}
\end{frame}

% DOMAIN C - QUESTION 29
\begin{frame}{Domain C: UNSCR 1373 vs 1267}
\scenario{A bank identifies a customer sending \$50,000 monthly in Bitcoin to a wallet. Not on UN sanctions list. Country has domestically designated recipient as terrorist financier under UNSCR 1373.}
\vspace{0.2cm}
\textbf{Must the bank freeze the assets?}
\vspace{0.2cm}
\begin{itemize}
    \item A) No; only UN 1267/1988 requires freezing
    \item B) Yes; UNSCR 1373 requires freezing domestically designated persons
    \item C) Only if amount exceeds \$100,000
    \item D) Only fiat assets must be frozen
\end{itemize}
\vspace{0.2cm}
\trap{No UN listing equals no obligation.}
\correct{Answer: B - UNSCR 1373 allows countries to create their own terrorist financing designation lists. Bank must comply with national law implementing 1373.}
\end{frame}

% DOMAIN C - QUESTION 30
\begin{frame}{Domain C: DNFBP Crypto Obligations}
\scenario{A law firm forms legal entities for crypto clients, provides registered agent services, holds client funds in escrow for crypto acquisitions, and facilitates bank introductions. Claims no AML obligations because not a VASP.}
\vspace{0.2cm}
\textbf{Does the law firm have AML obligations?}
\vspace{0.2cm}
\begin{itemize}
    \item A) No; law firms are exempt
    \item B) Yes; when acting as formation agent or holding client funds
    \item C) Only when directly handling crypto
    \item D) Only for sanctioned clients
\end{itemize}
\vspace{0.2cm}
\trap{Attorney-client privilege exempts.}
\correct{Answer: B - Under FATF Recommendation 23, DNFBPs including lawyers have AML obligations when organizing contributions for company creation, creating or operating legal persons, or managing client funds.}
\end{frame}

% DOMAIN D - QUESTION 31
\begin{frame}{Domain D: Blockchain Explorer Reliability}
\scenario{An investigator uses three blockchain explorers. Explorer A shows different transaction times than B. Explorer C misses 5 percent of transactions. Explorer A incorrectly labels an exchange as suspicious.}
\vspace{0.2cm}
\textbf{What is the best practice?}
\vspace{0.2cm}
\begin{itemize}
    \item A) Trust the most popular free explorer
    \item B) Compare data across multiple explorers; use paid tools for evidence
    \item C) Never use explorers
    \item D) Accept all explorer data as equally reliable
\end{itemize}
\vspace{0.2cm}
\trap{All explorers show same data.}
\correct{Answer: B - Blockchain explorers may present false information. Compare across multiple, verify reliability, use paid enterprise tools for court-admissible evidence.}
\end{frame}

% DOMAIN D - QUESTION 32
\begin{frame}{Domain D: Clustering Analysis}
\scenario{Address A and B are inputs to Transaction 1. Address B and C inputs to Transaction 2. Address C and D inputs to Transaction 3. Address D and E inputs to Transaction 4. Address E and A inputs to Transaction 5.}
\vspace{0.2cm}
\textbf{Which addresses likely belong to the same entity?}
\vspace{0.2cm}
\begin{itemize}
    \item A) No addresses can be grouped
    \item B) Addresses A, B, C, D, and E
    \item C) Only A and B
    \item D) Only D and E
\end{itemize}
\vspace{0.2cm}
\trap{Each address is independent.}
\correct{Answer: B - Common-input heuristic: multiple addresses as inputs to same transaction are likely controlled by same entity. Creates cluster containing A, B, C, D, E.}
\end{frame}

% DOMAIN D - QUESTION 33
\begin{frame}{Domain D: AI Transaction Monitoring}
\scenario{Legacy rules-based system: 10,000 alerts per month, 98 percent false positives, misses 30 percent of suspicious activity. AI system after tuning: 2,000 alerts per month, 60 percent false positives, 5 percent false negatives. But cannot explain decisions.}
\vspace{0.2cm}
\textbf{What are the trade-offs?}
\vspace{0.2cm}
\begin{itemize}
    \item A) AI is always better with no downsides
    \item B) AI reduces false positives and negatives but requires explainability
    \item C) Rules-based is always superior
    \item D) AI cannot be used for crypto
\end{itemize}
\vspace{0.2cm}
\trap{AI is perfect.}
\correct{Answer: B - AI improves detection but black-box models may not satisfy regulatory expectations. Requires explainability, ongoing testing, and governance.}
\end{frame}

% DOMAIN D - QUESTION 34
\begin{frame}{Domain D: Geolocation Red Flags}
\scenario{Customer withdraws \$500 cash from bank, then buys \$500 Bitcoin at ATM 30 minutes later, 50 miles from home. Pattern repeats 45 times. Income \$60,000; crypto purchases \$22,500 in 3 months.}
\vspace{0.2cm}
\textbf{What geolocation red flags exist?}
\vspace{0.2cm}
\begin{itemize}
    \item A) No red flags
    \item B) Distance from home, timing correlation, low ID thresholds, income inconsistency
    \item C) Only the withdrawal amount
    \item D) Only the income inconsistency
\end{itemize}
\vspace{0.2cm}
\trap{Travel distance is not suspicious.}
\correct{Answer: B - Geographic inconsistency, timing correlation between bank and crypto, low ID thresholds, and income inconsistency are multiple red flags requiring investigation.}
\end{frame}

% DOMAIN D - QUESTION 35
\begin{frame}{Domain D: Device Fingerprinting}
\scenario{Device fingerprints show: Device 1 accesses 50 accounts; Device 2 accesses 1 account from 200 IPs across 15 countries in 24 hours; Device 3 has failed logins; Device 4 known fraud device.}
\vspace{0.2cm}
\textbf{Which devices present highest risk?}
\vspace{0.2cm}
\begin{itemize}
    \item A) Device 4 only
    \item B) Devices 1, 2, 3, and 4
    \item C) Device 2 only
    \item D) Device 3 only
\end{itemize}
\vspace{0.2cm}
\trap{Only known fraud devices are risky.}
\correct{Answer: B - Device 1: mule network; Device 2: VPN/proxy abuse; Device 3: takeover attempt; Device 4: known fraud. All elevated risk.}
\end{frame}

% DOMAIN D - QUESTION 36
\begin{frame}{Domain D: Perpetual KYC}
\scenario{Traditional KYC: annual reviews. Perpetual KYC triggers reviews on: 200 percent volume increase, Grey List exposure, mixer usage, velocity changes, adverse media. Identifies 1,500 risks in first month missed by annual review.}
\vspace{0.2cm}
\textbf{What are advantages of perpetual KYC?}
\vspace{0.2cm}
\begin{itemize}
    \item A) No advantages; periodic KYC sufficient
    \item B) Real-time risk detection, immediate triggers, no undetected gaps
    \item C) Only reduces paperwork
    \item D) Only works for traditional banking
\end{itemize}
\vspace{0.2cm}
\trap{Annual KYC is sufficient.}
\correct{Answer: B - Perpetual KYC provides real-time detection of risk changes, immediate trigger-based reviews, eliminates window between reviews, and allocates resources to active risks.}
\end{frame}

% DOMAIN D - QUESTION 37
\begin{frame}{Domain D: OSINT for Crypto Investigations}
\scenario{Investigator traces wallet to forum post, to Twitter username, to real name, to LinkedIn, to employer, to corporate registry, to director, to adverse media (fraud conviction).}
\vspace{0.2cm}
\textbf{What OSINT techniques were used and what are limitations?}
\vspace{0.2cm}
\begin{itemize}
    \item A) OSINT is not reliable
    \item B) Blockchain explorer, search engine, social media, corporate registry; limitations include source reliability and potential false associations
    \item C) Only blockchain explorers useful
    \item D) OSINT always 100 percent accurate
\end{itemize}
\vspace{0.2cm}
\trap{OSINT is always accurate.}
\correct{Answer: B - OSINT techniques valuable but limitations include source reliability, outdated information, potential false associations. Requires cross-verification across multiple sources.}
\end{frame}

% DOMAIN D - QUESTION 38
\begin{frame}{Domain D: Entity Resolution for Crypto}
\scenario{Company A and Company B share address, phone number, same director (John Smith), same wallet for exchange funding. John Smith has fraud conviction in adverse media.}
\vspace{0.2cm}
\textbf{What does entity resolution reveal?}
\vspace{0.2cm}
\begin{itemize}
    \item A) Companies are unrelated
    \item B) Common beneficial ownership; consolidated risk assessment needed
    \item C) Only shared address relevant
    \item D) Only shared phone number matters
\end{itemize}
\vspace{0.2cm}
\trap{Separate legal entities are independent.}
\correct{Answer: B - Entity resolution links common address, phone, individual, and wallets. Related entities require consolidated risk assessment with aggregated exposure.}
\end{frame}

% DOMAIN D - QUESTION 39
\begin{frame}{Domain D: RPA for SAR Filing}
\scenario{Manual SAR filing: 55-85 minutes each, 500 SARs per month equals 458-708 hours. RPA implementation: auto-extracts data, populates 40 of 45 fields, generates draft narrative. Reduces time to 15-25 minutes (70 percent reduction).}
\vspace{0.2cm}
\textbf{What are benefits and risks of RPA?}
\vspace{0.2cm}
\begin{itemize}
    \item A) No benefits; RPA cannot be used
    \item B) Efficiency, reduced errors, faster filing; risks include automating flawed processes
    \item C) Only risks; no benefits
    \item D) RPA eliminates need for human review
\end{itemize}
\vspace{0.2cm}
\trap{RPA is perfect; no risks.}
\correct{Answer: B - Benefits: 70 percent time reduction, 95 percent error reduction. Risks: automating flawed processes, over-reliance, template narratives may miss unique red flags. Human review still critical.}
\end{frame}

% DOMAIN D - QUESTION 40
\begin{frame}{Domain D: Synthetic Identity Fraud in Crypto}
\scenario{A criminal creates 500 synthetic identities using AI-generated photos, fake government IDs, and fabricated credit histories. Each identity opens an account at a different VASP, purchases \$5,000 in crypto using stolen credit cards, then transfers funds to a central wallet. Total stolen: \$2.5 million. None of the VASPs detected the fraud because each transaction individually appeared legitimate.}
\vspace{0.2cm}
\textbf{What technology could help VASPs detect this synthetic identity scheme?}
\vspace{0.2cm}
\begin{itemize}
    \item A) Blockchain explorers only
    \item B) Device fingerprinting, network analysis, and cross-VASP data sharing
    \item C) Manual review only
    \item D) No technology can detect synthetic identities
\end{itemize}
\vspace{0.2cm}
\trap{Synthetic identities look real; cannot be detected.}
\correct{Answer: B - Device fingerprinting detects same device across accounts. Network analysis links wallets. Cross-VASP data sharing (314(b)) identifies pattern across institutions. Individual VASPs cannot see the full scheme alone.}
\redflag{Synthetic identity fraud requires cross-institutional information sharing and device intelligence to detect.}
\end{frame}

% ANSWER KEY SUMMARY
\begin{frame}{Answer Key Summary}
\begin{multicols}{2}
\small
\begin{itemize}
    \item Q1: \correct{B}
    \item Q2: \correct{C}
    \item Q3: \correct{B}
    \item Q4: \correct{B}
    \item Q5: \correct{B}
    \item Q6: \correct{B}
    \item Q7: \correct{C}
    \item Q8: \correct{B}
    \item Q9: \correct{B}
    \item Q10: \correct{B}
    \item Q11: \correct{B}
    \item Q12: \correct{B}
    \item Q13: \correct{C}
    \item Q14: \correct{B}
    \item Q15: \correct{B}
    \item Q16: \correct{B}
    \item Q17: \correct{B}
    \item Q18: \correct{B}
    \item Q19: \correct{B}
    \item Q20: \correct{B}
    \item Q21: \correct{B}
    \item Q22: \correct{B}
    \item Q23: \correct{B}
    \item Q24: \correct{B}
    \item Q25: \correct{B}
    \item Q26: \correct{B}
    \item Q27: \correct{B}
    \item Q28: \correct{B}
    \item Q29: \correct{B}
    \item Q30: \correct{B}
    \item Q31: \correct{B}
    \item Q32: \correct{B}
    \item Q33: \correct{B}
    \item Q34: \correct{B}
    \item Q35: \correct{B}
    \item Q36: \correct{B}
    \item Q37: \correct{B}
    \item Q38: \correct{B}
    \item Q39: \correct{B}
    \item Q40: \correct{B}
\end{itemize}
\end{multicols}
\vspace{0.3cm}
\centering
\greennote{40 Questions Covering All 4 Domains}
\end{frame}

% KEY CONCEPTS RECAP
\begin{frame}{Key Concepts Recap}
\small
\begin{itemize}
    \item \keyterm{Microstructuring:} Small crypto transactions to avoid detection
    \item \keyterm{Mixers/Tumblers:} CoinJoin, Wasabi, Tornado Cash - make tracing nearly impossible
    \item \keyterm{Privacy Coins:} Monero, Zcash - higher risk than transparent blockchains
    \item \keyterm{Stablecoins:} Ease of fiat conversion equals money laundering risk
    \item \keyterm{NFT Wash Trading:} Self-laundering through overpriced NFT sales
    \item \keyterm{DeFi VASP Status:} Control and influence test - governance tokens plus admin keys
    \item \keyterm{Travel Rule:} Originator and beneficiary info required for VASPs same as banks
    \item \keyterm{MiCA:} Single EU license, passporting rights, EUR 10,000 cash limit
    \item \keyterm{US AML Act:} Crypto exchanges equal MSBs with full BSA requirements
    \item \keyterm{FATF Recommendation 15:} New product risk assessment BEFORE launch
    \item \keyterm{Clustering:} Common-input heuristic groups addresses by entity
    \item \keyterm{Perpetual KYC:} Real-time trigger-based reviews for crypto
    \item \keyterm{Entity Resolution:} Links separate legal entities with common control
    \item \keyterm{Synthetic Identity:} Requires cross-VASP data sharing and device fingerprinting
\end{itemize}
\vspace{0.3cm}
\centering
\greennote{Good luck on your CAMS exam!}
\end{frame}

\end{document}